%----------------------------------------------------------------------------------------
% TWEEDEGRAADSONGELIJKHEDEN. Alle opgaven zijn eigen werk.
%
% Opbouw: Probleem, Start, leerstof, en op een nieuwe pagina de Opdrachten en Test jezelf (geen Koppel in de 2de graad).
% Labels van dit hoofdstuk beginnen met h4-.
% Pakketten: dit hoofdstuk staat in pakket 2-5; in pakket 2 is het een keuzehoofdstuk (K2).
% Verschillen: ongelijkheden met een parameter (4-5), een bewijsopdracht (5).
%----------------------------------------------------------------------------------------

\mychapter{4}{Tweedegraadsongelijkheden}

\begin{probleem}
\vspace{-1mm}\naast{Een vuurpijl wordt recht omhoog geschoten. Na $t$ seconden is hij ongeveer $h(t) = -5t^2 + 40t$ meter hoog. Om veilig te zijn, mag hij pas ontploffen als hij hoger is dan $60$ m. Tussen welke tijdstippen is hij zo hoog? En de vuurwerkmaker wil dat de pijl ontploft op minstens $75$ m hoogte: binnen welk tijdsinterval moet de lont dan opgebrand zijn? Met een vergelijking vind je de tijdstippen waarop de pijl precies $60$ m hoog is. Maar hier is een heel interval het antwoord.}{\begin{tikzpicture}[x=.38cm,y=.03cm,font=\scriptsize]
  \roosterleeg[$t$][$h$]{0}{1}{8}{0}{20}{80}
  \xgetallen{2,4,6,8}\ygetallen{20,40,60,80}
  \parabool{-5}{4}{80}{0}{8}
  \draw[hulplijn] (0,60) -- (8,60);
\end{tikzpicture}}
\end{probleem}

\lkr{Leerplan \pakketnaam: \alleenvoor{2}{keuzedoel K2 (een school mag dit hoofdstuk weglaten)}\alleenvoor{3}{LPD 25-30}\alleenvoor{4,5}{LPD 29-34} (dit hoofdstuk: het tekenschema van een tweedegraadsfunctie, tweedegraadsongelijkheden algebraïsch en grafisch oplossen, en problemen met een tweedegraadsongelijkheid\alleenvoor{4,5}{; ongelijkheden met een parameter}).\alleenvoor{3,4,5}{ Welk nummer precies bij dit hoofdstuk hoort, is na te kijken in het leerplan.} Voorkennis: hoofdstuk 1 tot 3, het tekenschema, de intervallen en de unie uit het handboek Functies. Het probleem komt terug in voorbeeld~\ref{vb:h4-terug}.}

In hoofdstuk 3 loste je vergelijkingen op: wanneer is de vuurpijl \textit{precies} $60$ m hoog? Vaak wil je iets anders weten: wanneer is hij \textit{hoger} dan $60$ m, wanneer is de winst \textit{positief}, voor welke breedte is de oppervlakte \textit{minstens} zo groot? Dan los je een ongelijkheid op. Het antwoord is geen lijstje getallen, maar een interval of een unie van intervallen. De sleutel is het tekenschema: waar ligt de parabool boven de $x$-as, en waar eronder?

\begin{oefeningen}[Start]
\oef{h4-start-1}

\oef{h4-verken}
\end{oefeningen}

%========================================================================================
\section{Het tekenschema}

In het handboek Functies maakte je het tekenschema van een functie: waar is $f(x)$ positief, waar negatief, en waar nul? Een parabool maakt geen sprongen. Haar teken kan dus alleen veranderen in een nulwaarde, waar de grafiek door de $x$-as gaat. Het tekenschema hangt dus af van de nulwaarden, en die ken je uit hoofdstuk 3. Er zijn drie gevallen: twee, één of geen nulwaarden.

\subsection{Twee nulwaarden}

Neem $f(x) = x^2 - 2x - 3 = (x + 1)(x - 3)$, met nulwaarden $-1$ en $3$. De parabool is naar boven geopend. Links van $-1$ en rechts van $3$ ligt ze boven de $x$-as, ertussen eronder. Bij $g(x) = -x^2 + 2x + 3$ is het net andersom: dezelfde nulwaarden, maar de parabool is naar onder geopend.

\figuren{\begin{tabular}{c@{\qquad\qquad}c}
\begin{tikzpicture}[x=.42cm,y=.3cm]
  \roosterleeg{-2}{1}{4}{-4}{1}{5}\xgetallen{-1,3}
  \fill[accent!15] plot[domain=-1:3,samples=40] (\x,{(\x+1)*(\x-3)}) -- cycle;
  \parabool{1}{1}{-4}{-1.95}{3.95}
  \punt{-1}{0}\punt{3}{0}
\end{tikzpicture} &
\begin{tikzpicture}[x=.42cm,y=.3cm]
  \roosterleeg{-2}{1}{4}{-4}{1}{5}\xgetallen{-1,3}
  \fill[accent!15] plot[domain=-1:3,samples=40] (\x,{-(\x+1)*(\x-3)}) -- cycle;
  \parabool{-1}{1}{4}{-1.95}{3.95}
  \punt{-1}{0}\punt{3}{0}
\end{tikzpicture} \\[1mm]
\tekenschema{5}{ & -1 & & 3 & }{+ & 0 & - & 0 & +} & \tekenschema[g(x)]{5}{ & -1 & & 3 & }{- & 0 & + & 0 & -}
\end{tabular}}

\begin{theorie}
\begin{stelling}{Tekenschema bij twee nulwaarden}
Heeft $f(x) = ax^2 + bx + c$ twee nulwaarden $x_1 < x_2$ ($D > 0$), dan heeft $f(x)$

\qquad \textbf{\textcolor{accent}{buiten}} de nulwaarden het teken van $\formule{a}$,

\qquad \textbf{\textcolor{accent}{tussen}} de nulwaarden het tegengestelde teken van $\formule{a}$.
\end{stelling}
\end{theorie}

Je hoeft de parabool dus niet te tekenen: de nulwaarden en het teken van $a$ zijn genoeg. Maar een kleine schets helpt om je niet te vergissen.

\voorbeeld{}{Maak het tekenschema van $f(x) = -2x^2 + 4x + 6$.\label{vb:h4-twee}}{%
\naast{\uitleg{We zoeken de nulwaarden, en kijken naar het teken van $a$.}
\begin{stappenvrij}{4.4cm}{4cm}
-2x^2 + 4x + 6 = 0 & \\
x^2 - 2x - 3 = 0 & Deel door $-2$. \\
(x + 1)(x - 3) = 0 & Nulwaarden $-1$ en $3$. \\
a = -2 < 0 & Buiten: $-$, tussen: $+$.
\end{stappenvrij}
\medskip
\tekenschema{5}{ & -1 & & 3 & }{- & 0 & + & 0 & -}}{\begin{tikzpicture}[x=.42cm,y=.2cm]
  \rooster{-2}{1}{4}{-4}{2}{8}
  \parabool{-2}{1}{8}{-1.5}{3.5}
  \punt{-1}{0}\punt{3}{0}
\end{tikzpicture}}}

\begin{oefeningen}
\oef{h4-factoren}
\end{oefeningen}

\subsection{Eén nulwaarde}

Is $D = 0$, dan raakt de parabool de $x$-as in haar top. Bijvoorbeeld $f(x) = x^2 - 6x + 9 = (x - 3)^2$. Een kwadraat is nooit negatief, dus $f(x)$ is overal positief, behalve in $3$, waar het nul is.

\begin{theorie}
\begin{stelling}{Tekenschema bij één nulwaarde}
Heeft $f(x) = ax^2 + bx + c$ één nulwaarde $x_1$ ($D = 0$), dan heeft $f(x)$ overal het teken van $\formule{a}$, behalve in $x_1$, waar $f(x_1) = 0$.
\end{stelling}
\end{theorie}

\voorbeeld{}{Maak het tekenschema van $f(x) = -x^2 + 4x - 4$.\label{vb:h4-een}}{%
\naast{\uitleg{We berekenen de discriminant.}
\begin{stappenvrij}{4.4cm}{4cm}
D = 16 - 4 \cdot (-1) \cdot (-4) = 0 & Eén nulwaarde. \\[1ex]
x = -\dfrac{4}{2 \cdot (-1)} = 2 & De top ligt op de $x$-as. \\[1ex]
a = -1 < 0 & Overal $-$, behalve in $2$.
\end{stappenvrij}
\medskip
\tekenschema{3}{ & 2 & }{- & 0 & -}}{\begin{tikzpicture}[x=.42cm,y=.3cm]
  \rooster{-1}{1}{5}{-5}{1}{2}
  \parabool{-1}{2}{0}{-.2}{4.2}
  \punt{2}{0}
\end{tikzpicture}}}

\subsection{Geen nulwaarden}

Is $D < 0$, dan snijdt de parabool de $x$-as niet. Ze ligt er volledig boven of volledig onder, en het teken van $a$ zegt welke van de twee.

\begin{theorie}
\begin{stelling}{Tekenschema zonder nulwaarden}
Heeft $f(x) = ax^2 + bx + c$ geen nulwaarden ($D < 0$), dan heeft $f(x)$ voor elke $x$ het teken van $\formule{a}$.
\end{stelling}
\end{theorie}

\voorbeeld{}{Maak het tekenschema van $f(x) = x^2 + 2x + 5$.\label{vb:h4-geen}}{%
\naast{\uitleg{We berekenen de discriminant.}
\begin{stappenvrij}{4.4cm}{4cm}
D = 4 - 20 = -16 < 0 & Geen nulwaarden. \\
a = 1 > 0 & Overal positief.
\end{stappenvrij}
\medskip
\tekenschema{1}{ }{+}

\medskip
\textit{Het tekenschema heeft geen nulwaarden: er staat alleen een $+$.}}{\begin{tikzpicture}[x=.4cm,y=.25cm]
  \rooster{-4}{1}{2}{-1}{1}{8}
  \parabool{1}{-1}{4}{-3}{1}
\end{tikzpicture}}}

\begin{werkwijze}
Bereken de nulwaarden (ontbinden of de discriminant). Zet ze van klein naar groot in het tekenschema, met een $0$ eronder. Vul het teken in: bij twee nulwaarden buiten het teken van $a$ en ertussen het tegengestelde; bij één of geen nulwaarde overal het teken van $a$. Controleer met één getal, bijvoorbeeld $x = 0$: $f(0) = c$.
\end{werkwijze}

\begin{oefeningen}
\oef{h4-teken}
\end{oefeningen}

%========================================================================================
\section{Ongelijkheden oplossen}

\subsection{Met het tekenschema}

Voor welke $x$ is $x^2 - x - 6 < 0$? Je zoekt de $x$ waarvoor de functie $f(x) = x^2 - x - 6$ negatief is. Dat staat in het tekenschema.

\begin{theorie}
Een \textbf{\textcolor{accent}{tweedegraadsongelijkheid}} is een ongelijkheid die je kunt herleiden tot

\qquad $\formule{ax^2 + bx + c > 0}$ \qquad met $a \neq 0$,

of dezelfde vorm met $<$, $\geqslant$ of $\leqslant$.
\end{theorie}

\voorbeeld{}{Los op: $x^2 - x - 6 < 0$.\label{vb:h4-ong}}{%
\naast{\uitleg{We maken het tekenschema van $x^2 - x - 6$ en zoeken waar het teken $-$ is.}
\begin{stappenvrij}{4.1cm}{4.3cm}
(x + 2)(x - 3) = 0 & Nulwaarden $-2$ en $3$. \\
a = 1 > 0 & Tussen de nulwaarden: $-$.
\end{stappenvrij}
\medskip
\tekenschema{5}{ & -2 & & 3 & }{+ & 0 & - & 0 & +}

\medskip
\begin{stappenvrij}{4.1cm}{4.3cm}
V = \ivoo{-2}{3} & De grenzen niet: daar is $f(x) = 0$, niet kleiner.
\end{stappenvrij}}{\begin{tikzpicture}[x=.42cm,y=.25cm]
  \rooster{-3}{1}{4}{-7}{1}{3}
  \asdeel{-2}{3}
  \parabool{1}{.5}{-6.25}{-2.6}{3.6}
  \openpunt{-2}{0}\openpunt{3}{0}
\end{tikzpicture}}}

Bij $\leqslant$ of $\geqslant$ horen de nulwaarden er wel bij: dan krijg je vierkante haken die naar het interval toe wijzen. Ligt de oplossing aan twee kanten, dan noteer je ze met de unie $\cup$, zoals in het handboek Functies.

\voorbeeld{}{Los op: $-2x^2 + 8x \leqslant 0$.\label{vb:h4-unie}}{%
\naast{\uitleg{De constante term ontbreekt: we brengen $x$ buiten haakjes.}
\begin{stappenvrij}{4.1cm}{4.3cm}
-2x(x - 4) = 0 & Nulwaarden $0$ en $4$. \\
a = -2 < 0 & Buiten de nulwaarden: $-$.
\end{stappenvrij}
\medskip
\tekenschema{5}{ & 0 & & 4 & }{- & 0 & + & 0 & -}

\medskip
\begin{stappenvrij}{4.1cm}{4.3cm}
V = \ivog{\minf}{0} \cup \ivgo{4}{\pinf} & Teken $-$ of $0$: de nulwaarden horen erbij.
\end{stappenvrij}}{\begin{tikzpicture}[x=.42cm,y=.2cm]
  \rooster{-2}{1}{6}{-6}{2}{10}
  \asdeel{-2}{0}\asdeel{4}{6}
  \parabool{-2}{2}{8}{-.65}{4.65}
  \punt{0}{0}\punt{4}{0}
\end{tikzpicture}}}

\begin{werkwijze}
Herleid de ongelijkheid tot $ax^2 + bx + c$ tegenover $0$. Maak het tekenschema van het linkerlid. Lees af waar het teken klopt met de ongelijkheid. Bij $<$ en $>$ horen de nulwaarden er niet bij, bij $\leqslant$ en $\geqslant$ wel. Noteer de oplossing met intervallen.
\end{werkwijze}

Een veelgemaakte fout is de vierkantswortel nemen. Uit $x^2 < 9$ volgt niet $x < 3$, want ook $x = -5$ voldoet aan $x < 3$, terwijl $(-5)^2 = 25$. Met het tekenschema van $x^2 - 9$ vind je de juiste oplossing: $\ivoo{-3}{3}$.

\begin{oefeningen}
\oef{h4-wortelfout}

\oef{h4-ongelijkheid}
\end{oefeningen}

\subsection{Eén of geen nulwaarden}

Heeft het linkerlid één of geen nulwaarden, dan is het tekenschema bijna overal hetzelfde teken. De oplossing is dan vaak heel $\R$, één getal, of helemaal niets. Het volgende voorbeeld draait om dat verschil, met drie ongelijkheden die op elkaar lijken.

\voorbeeld{}{Los op: $(x - 3)^2 > 0$, \ $(x - 3)^2 \geqslant 0$ \ en \ $(x - 3)^2 \leqslant 0$.\label{vb:h4-bijzonder}}{%
\naast{\uitleg{Er is één nulwaarde, $3$, en overal elders is het kwadraat positief.}
\begin{stappenvrij}{4.1cm}{4.3cm}
\multicolumn{2}{@{}l}{\tekenschema{3}{ & 3 & }{+ & 0 & +}} \\[2ex]
(x - 3)^2 > 0 & \\
V = \R \setminus \{3\} & Alles behalve $3$. \\
(x - 3)^2 \geqslant 0 & \\
V = \R & Ook $3$ voldoet. \\
(x - 3)^2 \leqslant 0 & \\
V = \{3\} & Alleen $3$: nergens negatief.
\end{stappenvrij}}{\begin{tikzpicture}[x=.42cm,y=.3cm]
  \rooster{-1}{1}{6}{-1}{1}{5}
  \parabool{1}{3}{0}{.8}{5.2}
  \punt{3}{0}
\end{tikzpicture}}}

Zo is ook $x^2 + 2x + 5 < 0$ zonder oplossing, want het linkerlid is overal positief (voorbeeld~\ref{vb:h4-geen}): $V = \emptyset$. En $x^2 + 2x + 5 > 0$ geldt voor elke $x$: $V = \R$.

\begin{oefeningen}
\oef{h4-bijzonder}
\end{oefeningen}

%========================================================================================
\section{Grafisch oplossen en toepassingen}

\subsection{Een parabool en een rechte}

Voor welke $x$ ligt de parabool $y = x^2$ boven de rechte $y = x + 2$? In hoofdstuk 3 vonden we de snijpunten: $(-1, 1)$ en $(2, 4)$. Op de grafiek zie je meteen waar de parabool boven de rechte ligt. Algebraïsch breng je alles naar één kant: $x^2 > x + 2$ wordt $x^2 - x - 2 > 0$.

\voorbeeld{}{Los op: $x^2 > x + 2$, grafisch en algebraïsch.\label{vb:h4-grafisch}}{%
\naast{\uitleg{Grafisch: waar ligt de parabool boven de rechte? Algebraïsch: herleid op nul en maak een tekenschema.}
\begin{stappenvrij}{4.1cm}{4.3cm}
x^2 - x - 2 > 0 & Alles naar het linkerlid. \\
(x - 2)(x + 1) = 0 & Nulwaarden $-1$ en $2$. \\
\multicolumn{2}{@{}l}{\tekenschema[x^2 - x - 2]{5}{ & -1 & & 2 & }{+ & 0 & - & 0 & +}} \\[2ex]
V = \ivoo{\minf}{-1} \cup \ivoo{2}{\pinf} & Links en rechts van de snijpunten.
\end{stappenvrij}}{\begin{tikzpicture}[x=.5cm,y=.36cm]
  \rooster{-3}{1}{3}{-1}{1}{6}
  \asdeel{-3}{-1}\asdeel{2}{3}
  \parabool{1}{0}{0}{-2.45}{2.45}
  \rechte[inkt]{-3}{3}{1}{2}
  \openpunt{2}{4}\openpunt{-1}{1}
  \aflezen{2}{4}\aflezen{-1}{1}
\end{tikzpicture}}}

\begin{oefeningen}
\oef{h4-grafisch}
\end{oefeningen}

\begin{alleenvoor}{4,5}
\subsection{Ongelijkheden met een parameter}

Een ongelijkheid duikt ook op als je wilt weten wanneer een vergelijking oplossingen heeft. Het aantal oplossingen hangt af van het teken van de discriminant. Staat er een onbekend getal in de vergelijking, dan wordt dat een ongelijkheid in dat getal.

\voorbeeld{}{Voor welke waarden van $m$ heeft $x^2 + mx + 9 = 0$ twee oplossingen?\label{vb:h4-parameter}}{%
\naast{\uitleg{Twee oplossingen betekent $D > 0$. De discriminant hangt af van $m$: we lossen een ongelijkheid in $m$ op.}
\begin{stappenvrij}{4.1cm}{4.3cm}
D = m^2 - 36 & $a = 1$, $b = m$, $c = 9$. \\
m^2 - 36 > 0 & Twee oplossingen. \\
(m + 6)(m - 6) = 0 & Nulwaarden $-6$ en $6$. \\
m < -6 \ \text{ of } \ m > 6 & Buiten de nulwaarden: $+$.
\end{stappenvrij}
\textit{Voor $m = \pm6$ is er één oplossing, voor $-6 < m < 6$ geen.}}{\begin{tikzpicture}[x=.18cm,y=.05cm]
  \roosterleeg[$m$][$D$]{-10}{2}{10}{-40}{10}{60}\xgetallen{-6,6}\ygetallen{-36,40}
  \asdeel{-10}{-6}\asdeel{6}{10}
  \parabool{1}{0}{-36}{-9.8}{9.8}
  \openpunt{-6}{0}\openpunt{6}{0}
\end{tikzpicture}}}

\begin{oefeningen}
\oef{h4-parameter}
\end{oefeningen}
\end{alleenvoor}

\subsection{Problemen met een ongelijkheid}

Woorden als \textit{hoger dan}, \textit{minstens}, \textit{hoogstens} of \textit{winst} wijzen op een ongelijkheid. Let op het domein: een tijd of een lengte is niet negatief.

\begin{oefeningen}
\oef{h4-context}
\end{oefeningen}

\voorbeeld{}{Terug naar het probleem. De vuurpijl is na $t$ seconden $h(t) = -5t^2 + 40t$ meter hoog. Wanneer is hij hoger dan $60$ m? En wanneer minstens $75$ m?\label{vb:h4-terug}}{%
\uitleg{We lossen twee ongelijkheden op. Om met kleinere getallen te werken, delen we door $-5$. Let op: dan keert het ongelijkheidsteken om.}
\naast{\begin{stappenvrij}{4.2cm}{4.1cm}
-5t^2 + 40t > 60 & Hoger dan $60$ m. \\
-5t^2 + 40t - 60 > 0 & Herleid op nul. \\
t^2 - 8t + 12 < 0 & Deel door $-5$: teken om. \\
(t - 2)(t - 6) < 0 & Nulwaarden $2$ en $6$. \\
2 < t < 6 & Tussen de nulwaarden: $-$. \\[1ex]
-5t^2 + 40t \geqslant 75 & Minstens $75$ m. \\
t^2 - 8t + 15 \leqslant 0 & Herleid en deel door $-5$. \\
(t - 3)(t - 5) \leqslant 0 & Nulwaarden $3$ en $5$. \\
3 \leqslant t \leqslant 5 & Grenzen inbegrepen.
\end{stappenvrij}}{\begin{tikzpicture}[x=.38cm,y=.03cm]
  \roosterleeg[$t$][$h$]{0}{1}{8}{0}{20}{80}
  \xgetallen{2,4,6,8}\ygetallen{20,40,60,80}
  \asdeel{2}{6}
  \parabool{-5}{4}{80}{0}{8}
  \draw[hulplijn] (0,60) -- (8,60);
  \draw[hulplijn] (0,75) -- (8,75);
  \openpunt{2}{60}\openpunt{6}{60}\punt{3}{75}\punt{5}{75}
\end{tikzpicture}}
De vuurpijl is hoger dan $60$ m van $2$ tot $6$ seconden na de lancering, dus gedurende $4$ seconden. Op minstens $75$ m hoogte is hij van $3$ tot $5$ seconden: de lont moet tussen $3$ en $5$ seconden opgebrand zijn.}

%========================================================================================
\opdrachtenkop

\begin{oefeningen}
\needspace{6\baselineskip}
\oef{h4-o1}

\needspace{6\baselineskip}
\oef{h4-o2}

\needspace{6\baselineskip}
\oef{h4-o3}

\needspace{6\baselineskip}
\oef{h4-o24}

\needspace{6\baselineskip}
\oef{h4-o5}

\needspace{6\baselineskip}
\oef{h4-o4}

\needspace{6\baselineskip}
\oef{h4-o6}

\needspace{6\baselineskip}
\oef{h4-o7}

\needspace{6\baselineskip}
\oef{h4-o11}

\needspace{6\baselineskip}
\oef{h4-o9}

\needspace{6\baselineskip}
\oef{h4-o8}

\needspace{6\baselineskip}
\oef{h4-o13}

\needspace{6\baselineskip}
\oef{h4-o10}

\needspace{6\baselineskip}
\oef{h4-o20}

\needspace{6\baselineskip}
\oef{h4-o12}

\needspace{6\baselineskip}
\oef{h4-o15}

\needspace{6\baselineskip}
\oef{h4-o23}

\needspace{6\baselineskip}
\oef{h4-o17}

\needspace{6\baselineskip}
\oef{h4-o14}

\needspace{6\baselineskip}
\oef{h4-o16}

\needspace{6\baselineskip}
\oef{h4-o18}

\needspace{6\baselineskip}
\oef{h4-o19}

\needspace{6\baselineskip}
\oef{h4-o21}

\needspace{6\baselineskip}
\oef{h4-o22}
\end{oefeningen}

%----------------------------------------------------------------------------------------
\needspace{14\baselineskip}
\opdrachtsectie{Test jezelf}

\begin{oefeningen}
\needspace{6\baselineskip}
\oef{h4-t1}

\needspace{6\baselineskip}
\oef{h4-t2}

\needspace{6\baselineskip}
\oef{h4-t3}

\needspace{6\baselineskip}
\oef{h4-t4}

\needspace{6\baselineskip}
\oef{h4-t5}

\needspace{6\baselineskip}
\oef{h4-t6}

\needspace{6\baselineskip}
\oef{h4-t7}

\needspace{6\baselineskip}
\oef{h4-t8}
\end{oefeningen}

\herhaalschema{%
  \fv{oef:h4-t1}{vb:h4-twee}{oef:h4-teken}%
  \fv{oef:h4-t2}{vb:h4-een}{oef:h4-o1}%
  \fv{oef:h4-t3}{vb:h4-ong}{oef:h4-ongelijkheid}%
  \fv{oef:h4-t4}{vb:h4-unie}{oef:h4-o3}%
  \fv{oef:h4-t5}{vb:h4-bijzonder}{oef:h4-bijzonder}%
  \fv{oef:h4-t6}{vb:h4-grafisch}{oef:h4-o7}%
  \fv{oef:h4-t7}{vb:h4-terug}{oef:h4-context}%
  \alleenvoor{4,5}{\fv{oef:h4-t8}{vb:h4-parameter}{oef:h4-parameter}}}

\ifzelfstudie
\section*{Maak je eigen samenvatting}
Sluit dit boek en schrijf in je schrift: hoe het tekenschema van een tweedegraadsfunctie eruitziet bij twee, één en geen nulwaarden, hoe je een tweedegraadsongelijkheid oplost, wanneer de grenzen erbij horen, wat er gebeurt met het ongelijkheidsteken als je deelt door een negatief getal, waarom je bij $x^2 < 9$ geen wortel mag nemen, en hoe je een ongelijkheid grafisch oplost\alleenvoor{4,5}{, en hoe je met de discriminant een ongelijkheid met een parameter opstelt}. Vergelijk daarna met de kaders.
\fi
